\chapter{Introduction}
\label{ch:introduction}

\section{What \Morendo\ is}

\Morendo\ is a production rule engine: you describe the \emph{facts} of a
situation and write \emph{rules} of the form ``if these facts are present,
do this''. The engine keeps every fact in a working memory, keeps track of
which rules currently match, and when asked to \emph{fire} runs the actions
of the matching rules. It is an implementation of the RETE algorithm, the
same family as CLIPS, Jess and Drools, and it speaks the CLIPS language, so
most rule files written for CLIPS or Jess load with little or no change.

\Morendo\ began as the ``morendo'' branch of the Jamocha project by Peter
Lin (itself descended from Sumatra), and it is meant for research, learning
and experimentation. Version \version\ is a modernized engine: Java~21, a
one-command build, current libraries, and the package root
\fn{org.morendo}.

\section{The CLIPS language in \Morendo}

\Morendo\ reads the constructs and functions a CLIPS programmer expects:
\fn{deftemplate}, \fn{defrule} with the \fn{not}, \fn{exists},
\fn{and}, \fn{or}, \fn{forall} and \fn{test} conditional elements,
predicate and return-value constraints, \fn{deffacts} with a CLIPS
\fn{reset}, \fn{deffunction} bodies of any length, \fn{defglobal},
\fn{defmodule} with focus and auto-focus, and an agenda with salience,
strategies, \fn{agenda}, \fn{refresh} and \fn{halt}. The action side has
\fn{if}, \fn{while}, \fn{loop-for-count}, \fn{foreach}, \fn{progn},
\fn{break} and \fn{return}, nested calls as slot values in \fn{assert},
\fn{modify} and \fn{duplicate}, the fact functions \fn{fact-slot-value}
and friends, \fn{format}, file routers with \fn{open}, \fn{close},
\fn{printout}, \fn{readline} and \fn{read}, the string, list, math and
type functions, \fn{gensym} and \fn{system}. What it does not read is
listed in Appendix~\ref{app:differences}.

\section{Features beyond plain CLIPS}

\begin{itemize}
\item \textbf{Java objects as facts.} Any Java bean can be declared as a
  template; its instances are asserted as \emph{shadow facts} and rules match
  their properties directly. Chapter~\ref{ch:java-objects}.
\item \textbf{Rule properties.} Rules can bypass the agenda
  (\fn{no-agenda}), carry effective and expiration dates, check the
  expiry of the facts they match (\fn{temporal-activation}), and carry a
  version and a chaining direction. A rule may also have
  \emph{modification actions}, run when a fact it matched is retracted.
  Chapter~\ref{ch:rule-properties}.
\item \textbf{Queries.} \fn{defquery} runs a parameterized query over
  working memory; \fn{defgraphquery} does the same over a graph of nodes and
  edges. Chapter~\ref{ch:queries}.
\item \textbf{MOLAP cubes.} \fn{defcube} declares a multidimensional cube
  fed by the facts that match its schema; rules query it with the
  \fn{cubequery} conditional element and react when its measures change.
  Chapter~\ref{ch:molap}.
\item \textbf{Temporal facts.} Facts and objects can carry effective and
  expiration times, and the time functions compare instants.
  Chapter~\ref{ch:temporal}.
\item \textbf{\fn{only} and \fn{multiple}.} Two conditional elements from
  Haley Enterprise that match when exactly one, or more than one, fact
  satisfies a pattern. Chapter~\ref{ch:language}.
\item \textbf{Rule cost analysis and profiling.} The engine counts the
  network nodes a rule needs and reports a relative cost, and times its
  own operations per engine. Chapter~\ref{ch:analysis}.
\item \textbf{Engines for servers.} The service module keeps a pool of
  engines per rule application, initialized from a JSON configuration,
  and hands one to each request thread. Chapter~\ref{ch:embedding}.
\end{itemize}

\section{What \Morendo\ leaves out}

\Morendo\ deliberately does not support \emph{ordered facts}, the
\fn{(color red)} style of fact without a template. Every fact belongs to a
template, either a \fn{deftemplate} or a declared Java class. Ordered facts
are convenient for quick prototypes, but rule sets built on them tend to
need rewriting once the data model settles; \Morendo\ makes you design the
model first. Appendix~\ref{app:differences} lists the other differences
from CLIPS.

\section{How the engine works, in one page}

A rule set is compiled into a network. Each template gets an
\emph{object type node}; each literal or predicate constraint in a rule
pattern gets an \emph{alpha node} hanging off it; each pair of patterns that
share a variable gets a \emph{join node}; and each rule ends in a
\emph{terminal node}. Alpha nodes that test the same thing are shared
between rules.

When a fact is asserted it enters the network at its template's node,
passes the alpha tests it satisfies, and is joined with the facts that
already passed the other patterns of each rule. When a complete match
reaches a terminal node an \emph{activation} is placed on the
\emph{agenda}. Retracting a fact withdraws the activations it took part in.
\fn{(fire)} takes activations off the agenda in strategy order (depth by
default, with rule salience taking precedence) and runs their actions, which
usually assert, modify or retract further facts and so create or withdraw
more activations, until the agenda is empty.

The important consequence for the user is that matching happens
\emph{incrementally, at assertion time}, not when \fn{fire} is called.
That is what makes RETE fast. A rule defined after its facts is matched
against them when it is compiled (Section~\ref{sec:rules-before-facts}),
so the order of rules and data is a matter of style, not correctness.
